\section{Evaluation}
\label{sec:evaluation}

\subsection{Attribution and method}

The evaluation was designed and executed by the project teammate, Zhang Hanyu, and is
recorded in \path{evaluation/results.md} together with the benchmark scripts
\path{scripts/benchmark-attention.ts} and \path{scripts/benchmark-postage.ts}, merged
into the implementation branch as pull request~\#2 (merge commit
\code{0330091})~\cite{tap-eval}. This section reproduces his results without alteration
and adds the author's interpretation in the light of the design. Where the two are mixed
the text says so.

The evaluation asks three questions, one per requirement group:

\begin{enumerate}[label=Q\arabic*.]
  \item Does coalescing with escalation-first rendering reduce the notification context
        injected into the receiver's prompt, and keep the escalation visible? (FR1, FR2)
  \item Does prepaid postage impose an economic limit on repeated delivery? (FR5, FR7)
  \item Do priority messages remain visible when ordinary notifications are limited by a
        quota? (FR9, FR10)
\end{enumerate}

The experiments were run on Windows against the branch as merged. The teammate reports
that the project built after minor local compatibility fixes to TypeScript typing and
build scripts, but that the OWS native package
\path{@silicon-intern/ows-core-win32-x64-msvc} was unavailable in that environment. This
constrained which mechanisms could be exercised end to end and is discussed under Q3 and
in Section~\ref{sec:limitations}.

\subsection{Experiment 1: notification flood and attention reduction}
\label{sec:eval-flood}

\paragraph{Set-up.}
Three peer agents send $N$ ordinary chat messages in round-robin, after which one
transfer request requiring approval arrives. The arrival order is the worst case for a
FIFO renderer: the escalation is last. Two pipelines are compared on identical event
streams: the \emph{legacy} pipeline (append-only queue, per-event rendering, hard cap of
twenty lines, opaque ``$N$ more omitted'' footer) and the \emph{new} pipeline (coalescing
queue, escalation-first plan, grouped summary). $N \in \{10, 50, 100, 500, 1000\}$.
Token cost is estimated with the project's existing approximation
$\lceil \text{characters}/4 \rceil$; the figures are therefore estimated tokens, not
tokenizer counts. ``Escalation visible'' records whether the transfer request appears as a
rendered line.

\paragraph{Results.}

\begin{table}[htbp]
\centering\small
\caption{Experiment 1: estimated tokens injected per prompt and escalation visibility
(teammate's results, reproduced verbatim).}
\label{tab:flood}
\begin{tabular}{@{}rrrrll@{}}
\toprule
\textbf{Messages} & \textbf{Legacy tokens} & \textbf{New tokens} & \textbf{Reduction} & \textbf{Legacy escal.} & \textbf{New escal.} \\
\midrule
10   & 256 &  97 & 62.1\% & Visible & Visible \\
50   & 482 &  98 & 79.7\% & Hidden  & Visible \\
100  & 482 &  98 & 79.7\% & Hidden  & Visible \\
500  & 482 & 100 & 79.3\% & Hidden  & Visible \\
1000 & 483 & 100 & 79.3\% & Hidden  & Visible \\
\bottomrule
\end{tabular}
\end{table}

\paragraph{Interpretation.}
Three observations follow directly from the numbers, and each traces to a specific design
decision.

First, the new pipeline's cost is essentially \emph{flat}: 97--100 estimated tokens from
10 to 1000 messages. This is the coalescing invariant of \S\ref{sec:design-coalescing} at
work---three peers produce three counted lines regardless of $N$, and the only growth is
the width of the \code{(xN)} suffix as the count gains digits. The legacy pipeline
saturates at 482--483 tokens from $N=50$ onward because its cap of twenty lines is
already full at 50 messages; the reduction therefore stabilises at about 79\% rather than
growing without bound, which is the honest way to read the table. At $N=10$ the legacy
renderer is under its cap, so the gap is smaller (62.1\%).

Second, the legacy pipeline \emph{hides the escalation} at every flood size of 50 and
above, exactly the failure that motivated FR2: twenty chat lines fill the cap and the
approval request falls into the footer. The new pipeline renders it in every case,
because the plan partitions escalations first. The evaluation thus confirms both halves
of the coalescing increment---less context, and the right context.

Third, the absolute numbers matter less than their shape. A saving of roughly 380
estimated tokens per prompt is modest in isolation; but the block is injected on
\emph{every} prompt while the notifications are pending, in a streaming host it is
injected on every wake, and without coalescing the block's size is set by the most
talkative peer rather than by the operator. Bounding it is what makes a per-peer
attention price meaningful at all (principle P2).

\paragraph{Reproducibility.}
The scenario is synthetic and deterministic (fixed peer names, templated message text,
fixed timestamps), so re-running \code{bun scripts/benchmark-attention.ts} on any
platform should reproduce Table~\ref{tab:flood} exactly. The author was unable to
re-execute it in the environment used to prepare this report and makes no claim of
independent replication.

\subsection{Experiment 2: prepaid postage and economic message control}
\label{sec:eval-postage}

\paragraph{Set-up.}
The experiment drives the implemented \code{FilePostageLedger} directly, in a temporary
data directory. A sender is issued a credit of 0.002~USDC; the receiver's standard price is
0.001~USDC per message. Three debits are attempted with sequence numbers 1, 2, 3.

\paragraph{Results.}

\begin{table}[htbp]
\centering\small
\caption{Experiment 2: stamp debits against a 0.002~USDC credit at 0.001~USDC per message
(teammate's results, reproduced verbatim).}
\label{tab:postage}
\begin{tabular}{@{}rrlr@{}}
\toprule
\textbf{Message} & \textbf{Cost} & \textbf{Result} & \textbf{Remaining credit} \\
\midrule
1 & 0.001 USDC & Accepted & 0.001 USDC \\
2 & 0.001 USDC & Accepted & 0 USDC \\
3 & 0.001 USDC & Rejected & 0 USDC \\
\bottomrule
\end{tabular}
\end{table}

The third debit returns the rejection reason \code{insufficient\_credit}.

\paragraph{Interpretation.}
The ledger admits exactly $\lfloor 0.002 / 0.001 \rfloor = 2$ standard attention units
and then refuses, which is the behaviour FR7 specifies: one payment, many spends, a hard
stop. The rejection is the same \code{PostageRejection} object that the service wraps in
a \code{-32050} on the wire (Listing~\ref{lst:reject}), carrying \code{remaining: "0"} and
\code{required: "0.001"}, so a sender's daemon can conclude mechanically that a top-up is
the remedy. The experiment also shows the accounting is exact in integer micro-units: the
remaining balance is reported as \code{0}, not as a floating-point residue.

What the experiment does not show should be stated plainly. It exercises the debit rule in
isolation: it does not include the on-chain payment, the \code{postage/topup}
round-trip, the certificate, or the transport. Those paths are covered by the mock E2E
phase~7 (Table~\ref{tab:e2e}), which runs in continuous integration but was not part of
this evaluation and has not yet been executed against live XMTP.

\subsection{Experiment 3: priority messages and notification quota}
\label{sec:eval-quota}

\paragraph{Intended set-up.}
A grant holder with \code{notificationsPerWeek} exhausted sends ordinary chatter and one
priority-stamped message; the receiver's drained batch is inspected for a single folded
summary line and an unfolded ``Priority message from'' escalation.

\paragraph{Outcome.}
A direct local execution was attempted but could not be completed: the benchmark requires
the daemon's runtime, which loads the OWS native module, and
\path{@silicon-intern/ows-core-win32-x64-msvc} was not available in the Windows
environment. The teammate therefore documented the experiment on the basis of (i)~the
implemented quota logic in \code{notification-quota.ts}, (ii)~the mock E2E phase~8
scenarios, which assert precisely the intended outcome on the loopback transport, and
(iii)~source-level review of the escalation-exclusion rule. He characterises the block as
an environment limitation rather than a logic failure, and the author agrees: the phase-8
scenarios and the nine unit tests in \code{notification-quota.test.ts} (including
``never folds escalations'' and ``fails open when the trust store or ledger cannot
answer'') pass in CI on Linux. Nevertheless, Q3 is answered by tests written by the
implementer, not by an independent measurement, and this report records it as such.

\subsection{Summary of findings}

\begin{description}[style=nextline,leftmargin=1.2em]
  \item[Finding 1 --- Coalescing bounds attention cost and protects escalations.]
        Estimated injected tokens fall by 62.1\% at $N=10$ and by 79.3--79.7\% at
        $N \ge 50$, and the new pipeline's cost is flat in $N$; the escalation is visible in
        all five new-pipeline runs and hidden in four of five legacy runs
        (Table~\ref{tab:flood}).
  \item[Finding 2 --- Postage creates an explicit economic boundary.]
        A 0.002~USDC credit at 0.001~USDC per message admits two messages and rejects the
        third with \code{insufficient\_credit} (Table~\ref{tab:postage}).
  \item[Finding 3 --- Priority messages are preserved by design, not yet by measurement.]
        The escalation-exclusion rule is implemented, unit-tested and E2E-tested on the
        loopback transport; a standalone local benchmark was blocked by a platform
        dependency.
\end{description}
